\chapter{Sistema de arquivos}
\label{CH:FS}
%         é processo ou thread de kernel?
% 
%         o texto sobre logging (e parte do texto sobre buffer) assume que o leitor
%         já conhece bastante sobre como inodes e diretórios funcionam,
%         mas eles são introduzidos depois.
% 
% 	é necessário decidir entre processador e CPU.
% 	
% 	certifique-se de usar buffer, não bloco 
% 
% 	Montagem

O propósito de um sistema de arquivos é organizar e armazenar dados. Sistemas de arquivos
tipicamente suportam compartilhamento de dados entre usuários e aplicações, bem como
\indextext{persistência}
para que os dados ainda estejam disponíveis após uma reinicialização.

O sistema de arquivos do xv6 oferece arquivos, diretórios e nomes de caminho semelhantes aos do Unix
(consulte o Capítulo~\ref{CH:UNIX}), e armazena seus dados em um disco virtio para
garantir a persistência. O sistema de arquivos enfrenta
vários desafios:
\begin{itemize}

\item O sistema de arquivos precisa de estruturas de dados no disco para representar a árvore
de diretórios e arquivos nomeados, para registrar a identidade dos
blocos que armazenam o conteúdo de cada arquivo, e para registrar quais áreas
do disco estão livres.
\item O sistema de arquivos deve suportar
\indextext{recuperação de falhas}.
Ou seja, se ocorrer uma falha (por exemplo, falta de energia), o sistema de arquivos deve
ainda funcionar corretamente após um reinício. O risco é que uma falha possa
interromper uma sequência de atualizações e deixar as estruturas de dados no disco inconsistentes
(por exemplo, um bloco que é usado em um arquivo e marcado como livre).
\item Diferentes processos podem operar no sistema de arquivos ao mesmo tempo,
portanto o código do sistema de arquivos deve coordenar essas operações para manter os invariantes.
\item O acesso ao disco é ordens de magnitude mais lento que o acesso
à memória, portanto o sistema de arquivos deve manter um cache na memória
dos blocos mais utilizados.

\end{itemize}

O restante deste capítulo explica como o xv6 lida com esses desafios.

%% 
%%  -------------------------------------------
%% 
\section{Visão geral}

A implementação do sistema de arquivos do xv6 é
organizada em sete camadas, conforme mostrado na
Figura~\ref{fig:fslayer}.
A camada de disco lê e escreve blocos em um disco rígido virtio.
A camada de cache de buffer armazena em cache os blocos do disco e sincroniza o acesso a eles,
garantindo que apenas um processo do kernel por vez possa modificar os
dados armazenados em um determinado bloco. A camada de logging permite que camadas superiores
agrupem atualizações de vários blocos em uma
\indextext{transação},
e garante que os blocos sejam atualizados atomicamente
em caso de falhas (isto é, ou todos são atualizados ou nenhum).
A camada de inode fornece arquivos individuais, cada um representado como um
\indextext{inode}
com um número-i único
e alguns blocos contendo os dados do arquivo. A camada de diretório
implementa cada diretório como um tipo especial de
inode cujo conteúdo é uma sequência de entradas de diretório, cada uma contendo o
nome de um arquivo e seu número-i.
A camada de nomes de caminho fornece
nomes de caminho hierárquicos como
\lstinline{/usr/rtm/xv6/fs.c},
e os resolve com busca recursiva.
A camada de descritor de arquivo abstrai muitos recursos Unix (por exemplo, pipes, dispositivos,
arquivos etc.) usando a interface do sistema de arquivos, simplificando a vida dos
programadores de aplicações.

\begin{figure}[t]
\center
\includegraphics[scale=0.5]{fig/fslayer.pdf}
\caption{Camadas do sistema de arquivos do xv6.}
\label{fig:fslayer}
\end{figure}

O hardware de disco tradicionalmente apresenta os dados no
disco como uma sequência numerada de 
\textit{blocos} 
de 512 bytes
\index{bloco}
(também chamados de
\textit{setores}): 
\index{setor}
o setor 0 são os primeiros 512 bytes, o setor 1 são os próximos, e assim por diante. O tamanho do bloco
que um sistema operacional usa para seu sistema de arquivos pode ser diferente do
tamanho do setor usado pelo disco, mas tipicamente o tamanho do bloco é um múltiplo do
tamanho do setor. O xv6 mantém cópias dos blocos que foram lidos para a memória
em objetos do tipo
\lstinline{struct buf}
\lineref{kernel/buf.h:/^struct.buf/}.
Os
dados armazenados nessa estrutura às vezes não estão sincronizados com o disco: podem ainda
não ter sido lidos do disco (o disco está processando mas ainda não retornou o conteúdo
do setor), ou podem ter sido atualizados pelo software
mas ainda não foram escritos no disco.

O sistema de arquivos precisa de um plano para onde armazenar os inodes e
blocos de conteúdo no disco.
Para isso, o xv6 divide o disco em várias
seções, conforme mostra a
Figura~\ref{fig:fslayout}.
O sistema de arquivos não usa
o bloco 0 (ele contém o setor de boot). O bloco 1 é chamado de 
\indextext{superbloco}; 
ele contém metadados sobre o sistema de arquivos (o tamanho do sistema de arquivos em blocos, o
número de blocos de dados, o número de inodes e o número de blocos no
log). Os blocos a partir do 2 armazenam o log. Após o log vêm os inodes, com vários inodes por bloco. Depois
disso vêm os blocos de bitmap que rastreiam quais blocos de dados estão em uso.
Os blocos restantes são blocos de dados; cada um está marcado
como livre no bloco de bitmap ou contém conteúdo de um arquivo ou diretório.
O superbloco é preenchido por um programa separado, chamado
\indexcode{mkfs},
que constrói um sistema de arquivos inicial.

O restante deste capítulo discute cada camada, começando com o
cache de buffer.
Fique atento às situações em que boas abstrações nas camadas inferiores
facilitam o projeto das superiores.
%% 
%%  -------------------------------------------
%% 
\section{Camada de cache de buffer}
\label{s:bcache}

A cache de buffer tem dois trabalhos: (1) sincronizar o acesso aos blocos de disco para garantir
que apenas uma cópia de um bloco esteja na memória e que apenas uma thread do kernel por vez
use essa cópia; (2) armazenar em cache blocos populares para que eles não precisem ser lidos novamente
do disco lento. O código está em
\lstinline{bio.c}.

A interface principal exportada pela cache de buffer consiste em
\indexcode{bread}
e
\indexcode{bwrite};
a primeira obtém um
\indextext{buf}
contendo uma cópia de um bloco que pode ser lido ou modificado em memória, e a
segunda escreve um buffer modificado no bloco apropriado no disco.
Uma thread do kernel deve liberar um buffer chamando
\indexcode{brelse}
quando terminar de usá-lo.
A cache de buffer usa um sleep-lock por buffer para garantir
que apenas uma thread por vez use cada buffer
(e, portanto, cada bloco de disco);
\lstinline{bread}
retorna um buffer bloqueado, e
\lstinline{brelse}
libera o lock.

Vamos voltar à cache de buffer.
A cache de buffer tem um número fixo de buffers para armazenar blocos de disco,
o que significa que se o sistema de arquivos solicitar um bloco que não está
já na cache, a cache de buffer deve reciclar um buffer que atualmente
contenha algum outro bloco. A cache de buffer recicla o
buffer menos recentemente usado para o novo bloco. A suposição é que
o buffer menos recentemente usado é o menos provável de ser usado
novamente em breve.

\begin{figure}[t]
\center
\includegraphics[scale=0.5]{fig/fslayout.pdf}
\caption{Estrutura do sistema de arquivos xv6.}
\label{fig:fslayout}
\end{figure}
%% 
%%  -------------------------------------------
%% 
\section{Código: Cache de buffer}

A cache de buffer é uma lista duplamente encadeada de buffers.
A função
\indexcode{binit},
chamada por
\indexcode{main}
\lineref{kernel/main.c:/binit/},
inicializa a lista com os
\indexcode{NBUF}
buffers no array estático
\lstinline{buf}
\linerefs{kernel/bio.c:/Create.linked.list/,/^..}/}.
Todo o restante do acesso à cache de buffer se refere à lista encadeada via
\indexcode{bcache.head},
não ao array
\lstinline{buf}.

Um buffer tem dois campos de estado associados a ele.
O campo
\indexcode{valid}
indica que o buffer contém uma cópia do bloco.
O campo \indexcode{disk}
indica que o conteúdo do buffer foi entregue ao
disco, o qual pode alterar o buffer (por exemplo, escrever
dados do disco em \lstinline{data}).

\lstinline{bread}
\lineref{kernel/bio.c:/^bread/}
chama
\indexcode{bget}
para obter um buffer para o setor fornecido
\lineref{kernel/bio.c:/b.=.bget/}.
Se o buffer precisar ser lido do disco,
\lstinline{bread}
chama
\indexcode{virtio_disk_rw}
para fazer isso antes de retornar o buffer.

\lstinline{bget}
\lineref{kernel/bio.c:/^bget/}
escaneia a lista de buffers por um buffer com os números de dispositivo e setor fornecidos
\linerefs{kernel/bio.c:/Is.the.block.already/,/^..}/}.
Se existir tal buffer,
\indexcode{bget}
adquire o sleep-lock do buffer.
\lstinline{bget}
então retorna o buffer bloqueado.

Se não houver buffer em cache para o setor fornecido,
\indexcode{bget}
deve criar um, possivelmente reutilizando um buffer que armazenava
um setor diferente.
Ele escaneia a lista de buffers uma segunda vez, procurando por um buffer
que não esteja em uso (\lstinline{b->refcnt = 0});
qualquer buffer assim pode ser usado.
\lstinline{bget}
edita os metadados do buffer para registrar o novo número de dispositivo e setor
e adquire seu sleep-lock.
Note que a atribuição
\lstinline{b->valid = 0}
garante que
\lstinline{bread}
irá ler os dados do bloco do disco
em vez de usar incorretamente o conteúdo anterior do buffer.

É importante que haja no máximo um buffer em cache por
setor de disco, para garantir que leitores vejam as escritas, e porque o
sistema de arquivos usa locks nos buffers para sincronização.
\lstinline{bget}
garante esse invariante mantendo o
\lstinline{bcache.lock}
continuamente desde a verificação no primeiro loop de se o
bloco está em cache até a declaração no segundo loop de que
o bloco agora está em cache (através da atribuição de
\lstinline{dev},
\lstinline{blockno},
e
\lstinline{refcnt}).
Isso faz com que a verificação da presença de um bloco e (se não
presente) a designação de um buffer para armazenar o bloco sejam atômicas.

É seguro para
\lstinline{bget}
adquirir o sleep-lock do buffer fora da
seção crítica protegida por \lstinline{bcache.lock},
já que o valor não-zero de
\lstinline{b->refcnt}
impede que o buffer seja reutilizado para um bloco de disco diferente.
O sleep-lock protege leituras
e escritas do conteúdo do buffer, enquanto o
\lstinline{bcache.lock}
protege informações sobre quais blocos estão em cache.

Se todos os buffers estiverem ocupados, então muitos processos estão
simultaneamente executando chamadas ao sistema de arquivos;
\lstinline{bget}
entra em pânico.
Uma resposta mais elegante poderia ser dormir até que um buffer ficasse livre,
embora isso traria a possibilidade de deadlock.

Uma vez que
\indexcode{bread}
tenha lido o disco (se necessário) e retornado o
buffer ao seu chamador, o chamador tem
uso exclusivo do buffer e pode ler ou escrever os bytes de dados.
Se o chamador modificar o buffer, ele deve chamar
\indexcode{bwrite}
para gravar os dados alterados no disco antes de liberar o buffer.
\lstinline{bwrite}
\lineref{kernel/bio.c:/^bwrite/}
chama
\indexcode{virtio_disk_rw}
para se comunicar com o hardware de disco.

Quando o chamador termina de usar um buffer,
ele deve chamar
\indexcode{brelse}
para liberá-lo.
(O nome
\lstinline{brelse},
uma abreviação de
b-release,
é críptico, mas vale a pena aprender:
ele se originou no Unix e é usado também no BSD, Linux e Solaris.)
\lstinline{brelse}
\lineref{kernel/bio.c:/^brelse/}
libera o sleep-lock e
move o buffer
para a frente da lista encadeada
\linerefs{kernel/bio.c:/b->next->prev.=.b->prev/,/bcache.head.next.=.b/}.
Mover o buffer faz com que a
lista fique ordenada por quão recentemente os buffers foram usados (ou seja, liberados):
o primeiro buffer na lista é o mais recentemente usado,
e o último é o menos recentemente usado.
Os dois loops em
\lstinline{bget}
aproveitam isso:
a varredura por um buffer existente pode ter que processar toda a lista
no pior caso, mas verificar os buffers mais recentemente usados
primeiro (começando em
\lstinline{bcache.head}
e seguindo os ponteiros
\lstinline{next})
reduz o tempo de varredura quando há boa localidade de referência.
A varredura para escolher um buffer a ser reutilizado escolhe o buffer menos recentemente usado
fazendo a varredura para trás
(seguindo os ponteiros
\lstinline{prev}).
%% 
%%  -------------------------------------------
%% 
\section{Camada de logging}

Um dos problemas mais interessantes no design de sistemas de arquivos é a recuperação após falhas. O problema surge porque muitas operações do sistema de arquivos envolvem múltiplas gravações no disco, e uma falha após apenas parte dessas gravações pode deixar o sistema de arquivos em disco em um estado inconsistente. Por exemplo, suponha que ocorra uma falha durante o truncamento de um arquivo (definindo seu tamanho como zero e liberando seus blocos de conteúdo). Dependendo da ordem das gravações no disco, a falha pode deixar um inode com uma referência a um bloco de conteúdo que está marcado como livre, ou pode deixar um bloco alocado, porém não referenciado.

O segundo caso é relativamente inofensivo, mas um inode que referencia um bloco já liberado pode causar sérios problemas após um reboot. Após reiniciar, o kernel pode alocar esse bloco para outro arquivo, e agora temos dois arquivos diferentes apontando, inadvertidamente, para o mesmo bloco. Se o xv6 suportasse múltiplos usuários, isso poderia ser um problema de segurança, já que o dono do arquivo antigo poderia ler e escrever nos blocos do novo arquivo, pertencente a outro usuário.

O xv6 resolve o problema de falhas durante operações no sistema de arquivos com uma forma simples de logging. Uma chamada de sistema no xv6 não grava diretamente nas estruturas de dados do sistema de arquivos no disco. Em vez disso, ela grava uma descrição de todas as operações de escrita que deseja realizar em um 
\indextext{log} 
no disco. Uma vez que a chamada de sistema tenha registrado todas as suas gravações, ela grava um registro especial de 
\indextext{commit}
no disco, indicando que o log contém uma operação completa. Nesse ponto, a chamada de sistema copia as gravações para as estruturas de dados reais no disco. Após essas gravações serem concluídas, o log é apagado do disco.

Se o sistema falhar e reiniciar, o código do sistema de arquivos realiza a recuperação da seguinte forma, antes de executar qualquer processo. Se o log estiver marcado como contendo uma operação completa, então o código de recuperação copia as gravações para seus devidos locais no sistema de arquivos. Se o log não estiver marcado como completo, o código de recuperação simplesmente o ignora. A recuperação finaliza apagando o log.

Por que o log do xv6 resolve o problema de falhas durante operações no sistema de arquivos? Se a falha ocorrer antes do commit, o log não estará marcado como completo, o código de recuperação o ignorará, e o estado do disco será como se a operação nunca tivesse iniciado. Se a falha ocorrer após o commit, então a recuperação irá repetir todas as gravações da operação, talvez reescrevendo blocos já parcialmente atualizados. Em ambos os casos, o log torna as operações atômicas em relação a falhas: após a recuperação, ou todas as gravações da operação estão no disco, ou nenhuma está.
%% 
%% 
%% 
\section{Design do log}

O log reside em um local fixo e conhecido, especificado no superbloco. Ele consiste em um bloco de cabeçalho seguido por uma sequência de cópias de blocos atualizados (``blocos logados''). O bloco de cabeçalho contém um array com números de setores, um para cada bloco logado, e a contagem de blocos no log. A contagem no bloco de cabeçalho no disco é zero, indicando que não há transação no log, ou diferente de zero, indicando que o log contém uma transação completa e confirmada com a quantidade indicada de blocos. O xv6 grava o bloco de cabeçalho quando uma transação é confirmada, mas não antes disso, e define a contagem como zero após copiar os blocos para o sistema de arquivos. Assim, uma falha no meio de uma transação resultará em uma contagem zero no cabeçalho; uma falha após o commit resultará em uma contagem diferente de zero.

O código de cada chamada de sistema indica o início e fim da sequência de gravações que devem ser atômicas em relação a falhas. Para permitir a execução concorrente de chamadas de sistema por diferentes processos, o sistema de logging pode acumular gravações de múltiplas chamadas em uma única transação. Assim, um único commit pode envolver as gravações de várias chamadas de sistema completas. Para evitar que uma chamada de sistema seja dividida entre transações, o sistema de logging apenas faz commit quando não há chamadas de sistema de arquivos em andamento.

A ideia de agrupar várias transações em um único commit é chamada de 
\indextext{group commit}.
Group commit reduz o número de operações no disco porque amortiza o custo fixo de um commit entre várias operações. Ele também entrega múltiplas gravações concorrentes ao sistema de disco ao mesmo tempo, possivelmente permitindo que o disco as grave todas em uma única rotação. O driver virtio do xv6 não suporta esse tipo de 
\indextext{batching},
mas o design do sistema de arquivos do xv6 o permite.

O xv6 dedica uma quantidade fixa de espaço no disco para armazenar o log. O número total de blocos gravados pelas chamadas de sistema em uma transação deve caber nesse espaço. Isso tem duas consequências. Nenhuma chamada de sistema pode gravar mais blocos distintos do que o espaço disponível no log. Isso não é problema para a maioria das chamadas, mas duas delas podem gravar muitos blocos:
\indexcode{write}
e
\indexcode{unlink}.
Uma gravação de arquivo grande pode gravar muitos blocos de dados, vários blocos de mapa de bits e um bloco de inode; remover um arquivo grande pode gravar vários blocos de mapa de bits e um inode. A chamada de sistema \lstinline{write} do xv6 divide grandes gravações em várias menores que cabem no log, e 
\lstinline{unlink}
não causa problemas porque na prática o sistema de arquivos do xv6 usa apenas um bloco de mapa de bits. A outra consequência do espaço limitado é que o sistema de logging não permite o início de uma nova chamada de sistema a menos que tenha certeza de que suas gravações caberão no espaço restante do log.
%% 
%% 
%% 
\section{Código: logging}

Um uso típico do log em uma chamada de sistema se parece com isso:
\begin{lstlisting}[]
  begin_op();
  ...
  bp = bread(...);
  bp->data[...] = ...;
  log_write(bp);
  ...
  end_op();
\end{lstlisting}

\indexcode{begin_op}
\lineref{kernel/log.c:/^begin.op/}
espera até que o sistema de logging não esteja realizando um commit, e até que haja espaço não reservado suficiente no log para as gravações desta chamada. 
\lstinline{log.outstanding}
conta o número de chamadas de sistema que reservaram espaço no log; o espaço total reservado é 
\lstinline{log.outstanding}
vezes
\lstinline{MAXOPBLOCKS}.
Incrementar
\lstinline{log.outstanding}
reserva espaço e impede que um commit ocorra durante essa chamada de sistema.
O código assume conservadoramente que cada chamada pode gravar até
\lstinline{MAXOPBLOCKS}
blocos distintos.

\indexcode{log_write}
\lineref{kernel/log.c:/^log.write/}
atua como um substituto para 
\indexcode{bwrite}.
Ele registra o número do setor do bloco em memória,
reservando um slot para ele no log do disco,
e fixa o buffer no cache de blocos para evitar que ele seja descartado.
O bloco deve permanecer no cache até o commit:
até lá, a cópia em cache é o único registro da modificação; ela não pode ser gravada diretamente em seu local no disco até o commit;
e outras leituras dentro da mesma transação devem ver as modificações.
\lstinline{log_write}
detecta quando um bloco é gravado várias vezes durante uma única transação, e aloca para ele o mesmo slot no log.
Essa otimização é comumente chamada de 
\indextext{absorção}.
É comum, por exemplo, que o bloco de disco contendo inodes de vários arquivos seja gravado diversas vezes dentro de uma transação. Absorvendo várias gravações em disco em uma só, o sistema de arquivos economiza espaço no log e melhora o desempenho, pois apenas uma cópia do bloco precisa ser gravada.

\indexcode{end_op}
\lineref{kernel/log.c:/^end.op/}
primeiro decrementa a contagem de chamadas de sistema pendentes.
Se essa contagem chega a zero, ele realiza o commit da transação atual chamando
\lstinline{commit().}
Esse processo ocorre em quatro etapas.
\lstinline{write_log()}
\lineref{kernel/log.c:/^write.log/}
copia cada bloco modificado da transação do cache para seu slot no log no disco.
\lstinline{write_head()}
\lineref{kernel/log.c:/^write.head/}
grava o bloco de cabeçalho no disco: esse é o ponto de commit, e uma falha após essa gravação fará com que a recuperação repita as gravações da transação a partir do log.
\indexcode{install_trans}
\lineref{kernel/log.c:/^install_trans/}
lê cada bloco do log e o grava em seu local apropriado no sistema de arquivos.
Finalmente, 
\lstinline{end_op}
grava o cabeçalho do log com contagem zero;
isso deve ocorrer antes da próxima transação começar a gravar blocos logados, para evitar que uma falha cause a recuperação usando o cabeçalho de uma transação com os blocos de outra.

\indexcode{recover_from_log}
\lineref{kernel/log.c:/^recover_from_log/}
é chamada de 
\indexcode{initlog}
\lineref{kernel/log.c:/^initlog/},
que por sua vez é chamada por \indexcode{fsinit}\lineref{kernel/fs.c:/^fsinit/} durante o boot, antes que o primeiro processo de usuário seja executado
\lineref{kernel/proc.c:/fsinit/}.
Ela lê o cabeçalho do log, e imita as ações de
\lstinline{end_op}
caso o cabeçalho indique que o log contém uma transação confirmada.

Um exemplo de uso do log ocorre em 
\indexcode{filewrite}
\lineref{kernel/file.c:/^filewrite/}.
A transação se parece com:
\begin{lstlisting}[]
      begin_op();
      ilock(f->ip);
      r = writei(f->ip, ...);
      iunlock(f->ip);
      end_op();
\end{lstlisting}
Esse código está envolto em um loop que divide grandes gravações em transações individuais de apenas alguns setores por vez, para evitar exceder o log. A chamada para
\indexcode{writei}
grava vários blocos como parte dessa transação: o inode do arquivo, um ou mais blocos de mapa de bits e alguns blocos de dados.
%% 
%% 
%% 
\section{Código: Alocador de blocos}

O conteúdo de arquivos e diretórios é armazenado em blocos de disco,
que devem ser alocados a partir de um conjunto livre.
O alocador de blocos do xv6
mantém um bitmap de blocos livres no disco, com um bit por bloco.
Um bit zero indica que o bloco correspondente está livre;
um bit um indica que está em uso.
O programa
\lstinline{mkfs}
define os bits correspondentes ao setor de boot, superbloco, blocos de log, blocos de inode
e blocos do bitmap.

O alocador de blocos fornece duas funções:
\indexcode{balloc}
aloca um novo bloco de disco, e
\indexcode{bfree}
libera um bloco.
O laço em
\lstinline{balloc}
na linha
\lineref{kernel/fs.c:/^..for.b.=.0/}
considera todos os blocos, começando no bloco 0 até
\lstinline{sb.size},
o número total de blocos no sistema de arquivos.
Ele procura por um bloco cujo bit no bitmap seja zero,
indicando que está livre.
Se
\lstinline{balloc}
encontra tal bloco, ele atualiza o bitmap
e retorna o bloco.
Por eficiência, o laço é dividido em duas partes.
O laço externo lê cada bloco de bits do bitmap.
O laço interno verifica todos os
Bits por Bloco (\lstinline{BPB})
em um único bloco de bitmap.
A condição de corrida que poderia ocorrer se dois processos tentassem alocar
um bloco ao mesmo tempo é evitada pelo fato de que
o cache de blocos permite apenas que um processo utilize um bloco de bitmap por vez.

\lstinline{bfree}
\lineref{kernel/fs.c:/^bfree/}
encontra o bloco de bitmap correto e limpa o bit correspondente.
Novamente, o uso exclusivo implícito por
\lstinline{bread}
e
\lstinline{brelse}
evita a necessidade de bloqueios explícitos.

Como ocorre com boa parte do código descrito no restante deste capítulo,
\lstinline{balloc}
e
\lstinline{bfree}
devem ser chamados dentro de uma transação.
%% 
%%  -------------------------------------------
%% 
\section{Camada de inodes}

O termo
\indextext{inode}
pode ter dois significados relacionados.
Pode se referir à estrutura de dados em disco que contém
o tamanho de um arquivo e a lista de números de blocos de dados.
Ou “inode” pode se referir a um inode em memória, que contém
uma cópia do inode em disco, bem como informações extras necessárias
dentro do kernel.

Os inodes em disco
são agrupados em uma área contígua
do disco chamada de blocos de inode.
Todos os inodes têm o mesmo tamanho, então é fácil, dado um
número n, encontrar o n-ésimo inode no disco.
Na verdade, esse número n, chamado de número do inode ou i-number,
é como os inodes são identificados na implementação.

O inode em disco é definido por uma
\indexcode{struct dinode}
\lineref{kernel/fs.h:/^struct.dinode/}.
O campo
\lstinline{type}
distingue entre arquivos, diretórios e arquivos especiais (dispositivos).
Um tipo zero indica que o inode em disco está livre.
O campo
\lstinline{nlink}
conta o número de entradas de diretório que
referenciam esse inode, a fim de reconhecer quando o
inode em disco e seus blocos de dados devem ser liberados.
O campo
\lstinline{size}
registra o número de bytes de conteúdo no arquivo.
O array
\lstinline{addrs}
armazena os números dos blocos de disco que contêm
o conteúdo do arquivo.

O kernel mantém o conjunto de inodes ativos em memória
em uma tabela chamada \indexcode{itable};
\indexcode{struct inode}
\lineref{kernel/file.h:/^struct.inode/}
é a cópia em memória de um
\lstinline{struct}
\lstinline{dinode}
em disco.
O kernel armazena um inode em memória apenas se houver
ponteiros C referenciando esse inode. O campo
\lstinline{ref}
conta o número de ponteiros C que referenciam o
inode em memória, e o kernel descarta o inode da
memória se essa contagem de referências cair para zero.
As funções
\indexcode{iget}
e
\indexcode{iput}
adquirem e liberam ponteiros para um inode,
modificando a contagem de referências.
Ponteiros para um inode podem vir de descritores de arquivo,
diretórios de trabalho atuais e código transitório do kernel,
como o
\lstinline{exec}.

Existem quatro mecanismos de bloqueio ou similares no código
de inode do xv6.
\lstinline{itable.lock}
protege o invariante de que um inode está presente na tabela de inodes
no máximo uma vez, e o invariante de que o campo
\lstinline{ref}
de um inode em memória conta o número de ponteiros em memória para o inode.
Cada inode em memória possui um campo
\lstinline{lock}
contendo um
sleep-lock, que garante acesso exclusivo aos campos do inode (como o tamanho do arquivo),
bem como aos blocos de conteúdo do arquivo ou diretório.
Um inode com
\lstinline{ref}
maior que zero é mantido na tabela,
e sua entrada na tabela não será reutilizada para outro inode.
Por fim, cada inode contém um campo
\lstinline{nlink}
(tanto em disco quanto em memória, se carregado) que
conta o número de entradas de diretório que referenciam o arquivo;
o xv6 não libera um inode se sua contagem de links for maior que zero.

Um ponteiro para
\lstinline{struct}
\lstinline{inode}
retornado por
\lstinline{iget()}
é garantido como válido até a chamada correspondente de
\lstinline{iput()};
o inode não será deletado, e a memória referenciada
pelo ponteiro não será reutilizada para outro inode.
\lstinline{iget()}
fornece acesso não exclusivo a um inode, de modo que
vários ponteiros podem referenciar o mesmo inode.
Muitas partes do código do sistema de arquivos dependem desse comportamento de
\lstinline{iget()},
tanto para manter referências de longo prazo a inodes (como arquivos abertos
e diretórios atuais), quanto para evitar condições de corrida
sem causar deadlock em código que manipula múltiplos inodes (como
a resolução de nomes de caminho).

A
\lstinline{struct}
\lstinline{inode}
retornada por
\lstinline{iget}
pode não conter informações úteis ainda.
Para garantir que ela contenha uma cópia do inode em disco,
o código deve chamar
\indexcode{ilock}.
Isso bloqueia o inode (impedindo que outro processo o
bloqueie com \lstinline{ilock}) e o lê do disco,
caso ainda não tenha sido lido.
\lstinline{iunlock}
libera o bloqueio do inode.
Separar a aquisição do ponteiro para o inode do bloqueio
ajuda a evitar deadlocks em certas situações, por exemplo durante
a busca em diretórios.
Vários processos podem manter um ponteiro C para um inode
retornado por
\lstinline{iget},
mas apenas um processo pode bloquear o inode por vez.

A tabela de inodes armazena apenas inodes para os quais o código do kernel
ou estruturas de dados mantêm ponteiros C.
Seu principal papel é sincronizar o acesso por múltiplos processos.
A tabela de inodes também funciona como cache para inodes frequentemente utilizados, mas
o cache é secundário; se um inode for usado com frequência, o cache de blocos provavelmente
o manterá em memória.
Código que modifica um inode em memória o escreve no disco com
\lstinline{iupdate}.

%% 
%%  -------------------------------------------
%% 
\section{Código: Inodes}

Para alocar um novo inode (por exemplo, ao criar um arquivo),
o xv6 chama
\indexcode{ialloc}
\lineref{kernel/fs.c:/^ialloc/}.
\lstinline{ialloc}
é similar a
\indexcode{balloc}:
ele percorre as estruturas de inode no disco, um bloco por vez,
procurando por uma que esteja marcada como livre.
Quando encontra uma, ele a reivindica escrevendo o novo 
\lstinline{type}
no disco e então retorna uma entrada da tabela de inodes
com a chamada de cauda para
\indexcode{iget}
\lineref{kernel/fs.c:/return.iget\(dev..inum\)/}.
O funcionamento correto de
\lstinline{ialloc}
depende do fato de que apenas um processo por vez
pode manter uma referência a
\lstinline{bp}:
\lstinline{ialloc}
pode ter certeza de que nenhum outro processo verá simultaneamente
que o inode está disponível e tentará reivindicá-lo.

\lstinline{iget}
\lineref{kernel/fs.c:/^iget/}
procura na tabela de inodes por uma entrada ativa (\lstinline{ip->ref}
\lstinline{>}
\lstinline{0})
com o dispositivo e número de inode desejados.
Se encontrar, retorna uma nova referência para esse inode
\linerefs{kernel/fs.c:/^....if.ip->ref.>.0/,/^....}/}.
Enquanto
\indexcode{iget}
varre a tabela, ele registra a posição do primeiro slot vazio
\linerefs{kernel/fs.c:/^....if.empty.==.0/,/empty.=.ip/},
que usará caso precise alocar uma nova entrada na tabela.

O código deve travar o inode usando
\indexcode{ilock}
antes de ler ou escrever seus metadados ou conteúdo.
\lstinline{ilock}
\lineref{kernel/fs.c:/^ilock/}
usa um sleep-lock para esse propósito.
Uma vez que
\indexcode{ilock}
tenha acesso exclusivo ao inode, ele o lê do disco (mais provavelmente do cache de buffer), se necessário.
A função
\indexcode{iunlock}
\lineref{kernel/fs.c:/^iunlock/}
libera o sleep-lock,
o que pode acordar processos que estavam esperando.

\lstinline{iput}
\lineref{kernel/fs.c:/^iput/}
libera um ponteiro C para um inode
decrementando o contador de referências
\lineref{kernel/fs.c:/^..ip->ref--/}.
Se esta for a última referência, o slot do inode na tabela de inodes
está agora livre e pode ser reutilizado para um inode diferente.

Se
\indexcode{iput}
perceber que não há mais referências C para um inode
e que o inode não tem mais links (ou seja, não está presente
em nenhum diretório), então o inode e seus blocos de dados devem
ser liberados.
\lstinline{iput}
chama
\indexcode{itrunc}
para truncar o arquivo a zero bytes, liberando os blocos de dados;
define o tipo do inode como 0 (não alocado);
e escreve o inode no disco
\lineref{kernel/fs.c:/inode.has.no.links.and/}.

O protocolo de travamento em
\indexcode{iput}
no caso em que ele libera o inode merece atenção especial.
Um perigo é que uma thread concorrente esteja esperando em
\lstinline{ilock}
para usar esse inode (por exemplo, para ler um arquivo ou listar um diretório),
e não esteja preparada para encontrar o inode não mais alocado.
Isso não pode acontecer porque não há como uma chamada de sistema
obter um ponteiro para um inode em memória se ele não tem links
e
\lstinline{ip->ref}
é um.
Essa única referência pertence à thread que está chamando
\lstinline{iput}.
O outro grande perigo é que uma chamada concorrente para
\lstinline{ialloc}
escolha o mesmo inode que
\lstinline{iput}
está liberando.
Isso só pode acontecer depois que
\lstinline{iupdate}
escreve no disco que o inode tem tipo zero.
Essa condição de corrida é benigna; a thread alocadora esperará
educadamente para adquirir o sleep-lock do inode antes de ler ou escrever,
momento no qual
\lstinline{iput}
já terminou com ele.

\lstinline{iput()}
pode escrever no disco. Isso significa que qualquer chamada de sistema que use o sistema de arquivos pode escrever no disco, pois pode ser a última a ter
uma referência ao arquivo. Mesmo chamadas como
\lstinline{read()}
que parecem apenas leitura, podem acabar chamando
\lstinline{iput().}
Isso, por sua vez, significa que até mesmo chamadas de sistema de leitura
devem ser encapsuladas em transações se utilizarem o sistema de arquivos.

Existe uma interação desafiadora entre
\lstinline{iput()}
e falhas.
\lstinline{iput()}
não trunca um arquivo imediatamente quando o número de links do arquivo
cai para zero, porque algum processo ainda pode manter uma referência ao inode
na memória: um processo ainda pode estar lendo ou escrevendo no arquivo, pois o abriu com sucesso.
Mas, se ocorrer uma falha antes que o último processo feche
o descritor de arquivo, o arquivo estará marcado como alocado no disco
mas nenhum diretório apontará para ele.

Sistemas de arquivos lidam com esse caso de duas formas. A solução simples é que,
após a recuperação (reboot), o sistema de arquivos escaneie todo o disco
procurando arquivos que estejam marcados como alocados, mas que não tenham nenhuma entrada de diretório apontando para eles.
Se encontrar tais arquivos, pode liberá-los.

A segunda solução não requer escanear o sistema de arquivos.
Neste caso, o sistema grava em disco (por exemplo, no superbloco) o número do inode
cujo número de links caiu para zero, mas cujo contador de referências ainda é diferente de zero.
Se o sistema de arquivos remover o arquivo quando o contador de referências chegar a 0,
então atualiza a lista em disco removendo aquele inode da lista.
Na recuperação, o sistema libera quaisquer arquivos que estiverem na lista.

O xv6 não implementa nenhuma das duas soluções, o que significa que inodes podem permanecer marcados como alocados
no disco mesmo que não estejam mais em uso.
Com o tempo, o xv6 corre o risco de ficar sem espaço em disco.

%% 
%% 
%% 
\section{Código: Conteúdo do Inode}

\begin{figure}[t]
\center
\includegraphics[scale=0.5]{fig/inode.pdf}
\caption{A representação de um arquivo no disco.}
\label{fig:inode}
\end{figure}

A estrutura de inode no disco,
\indexcode{struct dinode},
contém um tamanho e um vetor de números de bloco (veja
Figura~\ref{fig:inode}).
Os dados do inode são encontrados nos blocos listados
no vetor
\lstinline{addrs}
do
\lstinline{dinode}.
Os primeiros
\indexcode{NDIRECT}
blocos de dados estão listados nas primeiras
\lstinline{NDIRECT}
entradas do vetor; esses blocos são chamados de
\indextext{blocos diretos}.
Os próximos
\indexcode{NINDIRECT}
blocos de dados não estão listados no inode,
mas sim em um bloco de dados chamado
\indextext{bloco indireto}.
A última entrada do vetor
\lstinline{addrs}
guarda o endereço do bloco indireto.
Assim, os primeiros 12 kB (
\lstinline{NDIRECT} 
\lstinline{x}
\indexcode{BSIZE})
de um arquivo podem ser carregados diretamente dos blocos listados no inode,
enquanto os próximos
\lstinline{256} kB (
\lstinline{NINDIRECT}
\lstinline{x}
\lstinline{BSIZE})
só podem ser acessados consultando o bloco indireto.
Essa é uma boa representação em disco, porém
complexa para os clientes.
A função
\indexcode{bmap}
gerencia essa representação para que rotinas de mais alto nível, como
\indexcode{readi}
e
\indexcode{writei},
que veremos em breve, não precisem lidar com essa complexidade.
\lstinline{bmap}
retorna o número do bloco no disco do
\lstinline{bn}'ésimo
bloco de dados para o inode
\lstinline{ip}.
Se
\lstinline{ip}
ainda não possui esse bloco,
\lstinline{bmap}
o aloca.

A função
\indexcode{bmap}
\lineref{kernel/fs.c:/^bmap/}
começa tratando o caso simples: os primeiros
\indexcode{NDIRECT}
blocos estão listados diretamente no inode
\linerefs{kernel/fs.c:/^..if.bn.<.NDIRECT/,/^..}/}.
Os próximos
\indexcode{NINDIRECT}
blocos estão listados no bloco indireto em
\lstinline{ip->addrs[NDIRECT]}.
\lstinline{bmap}
lê o bloco indireto
\lineref{kernel/fs.c:/bp.=.bread.ip->dev..addr/}
e então lê o número do bloco na posição correta dentro do bloco
\lineref{kernel/fs.c:/a.=..uint\*.bp->data/}.
Se o número de bloco exceder
\lstinline{NDIRECT+NINDIRECT},
\lstinline{bmap} 
entra em pânico;
\lstinline{writei}
contém a verificação que evita que isso aconteça
\lineref{kernel/fs.c:/off...n...MAXFILE.BSIZE/}.

\lstinline{bmap}
aloca blocos conforme necessário.
Uma entrada zero em
\lstinline{ip->addrs[]} ou no bloco indireto
indica que nenhum bloco foi alocado.
À medida que
\lstinline{bmap}
encontra zeros, os substitui por números de blocos novos,
alocados sob demanda
\linerefs{kernel/fs.c:/^....if..addr.=.*==.0/,/./}
\linerefs{kernel/fs.c:/^....if..addr.*NDIRECT.*==.0/,/./}.

\indexcode{itrunc}
libera os blocos de um arquivo, redefinindo seu tamanho para zero.
\lstinline{itrunc}
\lineref{kernel/fs.c:/^itrunc/}
começa liberando os blocos diretos
\linerefs{kernel/fs.c:/^..for.i.=.0.*NDIRECT/,/^..}/},
depois os listados no bloco indireto
\linerefs{kernel/fs.c:/^....for.j.=.0.*NINDIRECT/,/^....}/},
e por fim o próprio bloco indireto
\linerefs{kernel/fs.c:/^....bfree.*NDIRECT/,/./}.

\lstinline{bmap}
facilita o acesso aos dados de um inode para
\indexcode{readi}
e
\indexcode{writei}.
\lstinline{readi}
\lineref{kernel/fs.c:/^readi/}
começa
garantindo que o deslocamento e o número de bytes não estejam
além do fim do arquivo.
Leituras que começam além do final retornam erro
\linerefs{kernel/fs.c:/^..if.off.>.ip->size/,/./}
enquanto leituras que começam ou cruzam o final do arquivo
retornam menos bytes do que o solicitado
\linerefs{kernel/fs.c:/^..if.off.\+.n.>.ip->size/,/./}.
O laço principal processa cada bloco do arquivo,
copiando dados do buffer para
\lstinline{dst}
\linerefs{kernel/fs.c:/^..for.tot=0/,/^..}/}.

\indexcode{writei}
\lineref{kernel/fs.c:/^writei/}
é idêntica à
\indexcode{readi},
com três exceções:
escritas que começam ou cruzam o fim do arquivo
fazem o arquivo crescer, até o tamanho máximo permitido
\linerefs{kernel/fs.c:/^..if.off.\+.n.>.MAXFILE/,/./};
o laço copia dados para os buffers em vez de lê-los
\lineref{kernel/fs.c:/either.copyin.*bp->data/};
e se a escrita expandiu o arquivo,
\indexcode{writei}
deve atualizar seu tamanho
\linerefs{kernel/fs.c:/^..if.off.>.ip->size\)/,/./}.

A função
\indexcode{stati}
\lineref{kernel/fs.c:/^stati\(/}
copia os metadados do inode para a estrutura
\lstinline{stat},
que é exposta aos programas de usuário
por meio da chamada de sistema
\indexcode{stat}.
%% 
%% 
%% 
\section{Código: camada de diretório}

Um diretório é implementado internamente de forma muito semelhante a um arquivo.
Seu inode tem tipo
\indexcode{T_DIR}
e seus dados são uma sequência de entradas de diretório.
Cada entrada é uma
\indexcode{struct dirent}
\lineref{kernel/fs.h:/^struct.dirent/},
que contém um nome e um número de inode.
O nome tem no máximo
\indexcode{DIRSIZ}
(14) caracteres;
se for menor, é terminado por um byte NULL (0).
Entradas de diretório com número de inode igual a zero estão livres.

A função
\indexcode{dirlookup}
\lineref{kernel/fs.c:/^dirlookup/}
procura uma entrada com o nome dado dentro de um diretório.
Se encontrar, retorna um ponteiro para o inode correspondente, destravado,
e define 
\lstinline{*poff}
como o deslocamento em bytes da entrada dentro do diretório,
caso o chamador deseje editá-la.
Se
\lstinline{dirlookup}
encontra uma entrada com o nome correto,
ela atualiza
\lstinline{*poff}
e retorna um inode destravado
obtido via
\indexcode{iget}.
\lstinline{dirlookup}
é o motivo pelo qual 
\lstinline{iget}
retorna inodes destravados.
O chamador trancou
\lstinline{dp},
então se a busca foi por
\indexcode{.},
um alias para o diretório atual,
tentar trancar o inode antes de retorná-lo tentaria re-trancar
\lstinline{dp}
e causaria um deadlock.
(Há cenários de deadlock mais complicados envolvendo
múltiplos processos e
\indexcode{..},
um alias para o diretório pai;
\lstinline{.}
não é o único problema.)
O chamador pode destrancar
\lstinline{dp}
e então trancar
\lstinline{ip},
garantindo que mantenha apenas um lock por vez.

A função
\indexcode{dirlink}
\lineref{kernel/fs.c:/^dirlink/}
escreve uma nova entrada de diretório com o nome e número de inode fornecidos no
diretório
\lstinline{dp}.
Se o nome já existir,
\lstinline{dirlink}
retorna um erro
\linerefs{kernel/fs.c:/Check.that.name.is.not.present/,/^..}/}.
O laço principal lê as entradas do diretório procurando uma entrada não alocada.
Quando encontra, interrompe o laço mais cedo
\linerefs{kernel/fs.c:/Look for an empty dirent/,/.*break/},
com 
\lstinline{off}
ajustado para o deslocamento da entrada disponível.
Caso contrário, o laço termina com
\lstinline{off}
ajustado para
\lstinline{dp->size}.
De qualquer forma, 
\lstinline{dirlink}
então adiciona uma nova entrada ao diretório
escrevendo no deslocamento
\lstinline{off}
\linerefs{kernel/fs.c:/if.writei/,/return -1/}.
%% 
%% 
%% 
\section{Código: Nomes de caminho}

A busca por nome de caminho envolve uma sucessão de chamadas a
\indexcode{dirlookup},
uma para cada componente do caminho.
\lstinline{namei}
\lineref{kernel/fs.c:/^namei/}
avalia 
\lstinline{path}
e retorna o 
\lstinline{inode}
correspondente.
A função
\indexcode{nameiparent}
é uma variante: ela para antes do último elemento, retornando o 
inode do diretório pai e copiando o último elemento em
\lstinline{name}.
Ambas chamam a função generalizada
\indexcode{namex}
para realizar o trabalho real.

\lstinline{namex}
\lineref{kernel/fs.c:/^namex/}
começa decidindo onde a avaliação do caminho se inicia.
Se o caminho começar com uma barra, a avaliação começa na raiz;
caso contrário, no diretório atual
\linerefs{kernel/fs.c:/..if.\*path.==....\)/,/idup/}.
Em seguida, usa
\indexcode{skipelem}
para considerar cada elemento do caminho por vez
\lineref{kernel/fs.c:/while.*skipelem/}.
Cada iteração do laço deve procurar 
\lstinline{name}
no inode atual
\lstinline{ip}.
A iteração começa trancando
\lstinline{ip}
e verificando se ele é um diretório.
Caso contrário, a busca falha
\linerefs{kernel/fs.c:/^....ilock.ip/,/^....}/}.
(Trancar
\lstinline{ip}
é necessário não porque 
\lstinline{ip->type}
possa mudar—não pode—mas porque
até que 
\indexcode{ilock}
rode,
\lstinline{ip->type}
não é garantido de ter sido carregado do disco.)
Se a chamada for 
\indexcode{nameiparent}
e este for o último elemento do caminho, o laço para mais cedo,
como previsto na definição de
\lstinline{nameiparent};
o último elemento do caminho já foi copiado para
\lstinline{name},
então
\indexcode{namex}
apenas precisa
retornar o
\lstinline{ip}
destravado
\linerefs{kernel/fs.c:/^....if.nameiparent/,/^....}/}.
Por fim, o laço procura o elemento do caminho usando
\indexcode{dirlookup}
e se prepara para a próxima iteração atribuindo
\lstinline{ip = next}
\linerefs{kernel/fs.c:/^....if..next.*dirlookup/,/^....ip.=.next/}.
Quando o laço fica sem elementos, ele retorna
\lstinline{ip}.

O procedimento
\lstinline{namex}
pode levar bastante tempo para concluir: ele pode envolver diversas operações de disco para
ler inodes e blocos de diretório para os diretórios percorridos no caminho
(se eles não estiverem no cache de buffer). O Xv6 foi cuidadosamente projetado para que, se uma
invocação de
\lstinline{namex}
por uma thread do kernel estiver bloqueada esperando E/S de disco, outra thread do kernel procurando por
um caminho diferente possa prosseguir paralelamente.
\lstinline{namex}
tranca cada diretório do caminho separadamente, permitindo que buscas em diretórios diferentes
ocorram em paralelo.

Essa concorrência introduz alguns desafios. Por exemplo, enquanto uma thread do kernel
procura um caminho, outra pode estar modificando a árvore de diretórios removendo um diretório.
Um risco potencial é que uma busca esteja acessando um diretório que foi excluído por outra thread e
seus blocos tenham sido reutilizados por outro diretório ou arquivo.

O Xv6 evita essas condições de corrida. Por exemplo, ao executar
\lstinline{dirlookup}
dentro de
\lstinline{namex},
a thread de busca mantém o lock no diretório e
\lstinline{dirlookup}
retorna um inode obtido usando
\lstinline{iget}.
\lstinline{iget}
aumenta o contador de referências do inode. Só depois de receber o
inode de
\lstinline{dirlookup}
é que
\lstinline{namex}
libera o lock do diretório. A partir desse ponto outra thread pode remover o inode do
diretório, mas o xv6 não o apagará ainda, porque o contador de referências
do inode ainda é maior que zero.

Outro risco é deadlock. Por exemplo,
\lstinline{next}
aponta para o mesmo inode que
\lstinline{ip}
ao procurar por ".".
Trancar
\lstinline{next}
antes de liberar o lock de
\lstinline{ip}
causaria um deadlock.
Para evitar isso,
\lstinline{namex}
destranca o diretório antes de obter o lock de
\lstinline{next}.
Aqui vemos novamente por que a separação entre
\lstinline{iget}
e
\lstinline{ilock}
é importante.
%% 
%% 
%% 
\section{Camada de descritor de arquivo}

Um aspecto interessante da interface Unix é que a maioria dos recursos do sistema
são representados como arquivos, incluindo dispositivos como o console, pipes e,
claro, arquivos reais. A camada de descritor de arquivos é responsável por
essa uniformidade.

O Xv6 dá a cada processo sua própria tabela de arquivos abertos, ou
descritores de arquivos, como vimos no
Capítulo~\ref{CH:UNIX}.
Cada arquivo aberto é representado por uma
\indexcode{struct file}
\lineref{kernel/file.h:/^struct.file/},
que encapsula um inode ou um pipe,
além de um deslocamento de E/S.
Cada chamada a 
\indexcode{open}
cria um novo arquivo aberto (uma nova
\lstinline{struct}
\lstinline{file}):
se múltiplos processos abrirem o mesmo arquivo de forma independente,
as instâncias terão deslocamentos de E/S diferentes.
Por outro lado, um único arquivo aberto
(a mesma
\lstinline{struct}
\lstinline{file})
pode aparecer
várias vezes na tabela de arquivos de um processo
e também nas tabelas de outros processos.
Isso acontece se um processo usa
\lstinline{open}
para abrir o arquivo e cria aliases usando
\indexcode{dup}
ou o compartilha com um filho via
\indexcode{fork}.
Um contador de referências acompanha o número de referências a
um determinado arquivo aberto.
Um arquivo pode estar aberto para leitura, escrita ou ambos.
Os campos
\lstinline{readable}
e
\lstinline{writable}
controlam isso.

Todos os arquivos abertos no sistema são mantidos em uma tabela global de arquivos,
a 
\indexcode{ftable}.
A tabela de arquivos
possui funções para alocar um arquivo
(\indexcode{filealloc}),
criar uma referência duplicada
(\indexcode{filedup}),
liberar uma referência
(\indexcode{fileclose}),
e ler e escrever dados
(\indexcode{fileread}
e 
\indexcode{filewrite}).

As três primeiras seguem um formato já familiar.
\lstinline{filealloc}
\lineref{kernel/file.c:/^filealloc/}
percorre a tabela de arquivos procurando por um arquivo não referenciado
(\lstinline{f->ref}
\lstinline{==}
\lstinline{0})
e retorna uma nova referência;
\indexcode{filedup}
\lineref{kernel/file.c:/^filedup/}
incrementa o contador de referências;
e
\indexcode{fileclose}
\lineref{kernel/file.c:/^fileclose/}
o decrementa.
Quando o contador de referências chega a zero,
\lstinline{fileclose}
libera o pipe ou inode subjacente,
conforme o tipo.

As funções
\indexcode{filestat},
\indexcode{fileread},
e
\indexcode{filewrite}
implementam as operações 
\indexcode{stat},
\indexcode{read},
e
\indexcode{write}
em arquivos.
\lstinline{filestat}
\lineref{kernel/file.c:/^filestat/}
só é permitida em inodes e chama
\indexcode{stati}.
\lstinline{fileread}
e
\lstinline{filewrite}
verificam se a operação é permitida pelo
modo de abertura e então
encaminham a chamada para a implementação correspondente,
seja pipe ou inode.
Se o arquivo representar um inode,
\lstinline{fileread}
e
\lstinline{filewrite}
usam o deslocamento de E/S como o deslocamento da operação
e depois o atualizam
\linerefs{kernel/file.c:/readi/,/./}
\linerefs{kernel/file.c:/writei/,/./}.
Pipes não têm conceito de deslocamento.
Lembre que as funções de inode exigem que o chamador
gerencie o locking
\linerefs{kernel/file.c:/stati/-1,/iunlock/}
\linerefs{kernel/file.c:/readi/-1,/iunlock/}
\linerefs{kernel/file.c:/writei\(f/-1,/iunlock/}.
O locking no inode tem o efeito colateral conveniente de que os
deslocamentos de leitura e escrita são atualizados atomicamente, de forma que
múltiplas escritas simultâneas no mesmo arquivo
não sobrescrevem os dados umas das outras, embora suas escritas possam acabar entrelaçadas.
%% 
%% 
%% 
\section{Código: Chamadas de sistema}

Com as funções fornecidas pelas camadas inferiores, a implementação da maioria
das chamadas de sistema é trivial
(veja
\fileref{kernel/sysfile.c}).
Há algumas chamadas que
merecem uma análise mais cuidadosa.

As funções
\indexcode{sys_link}
e
\indexcode{sys_unlink}
editam diretórios, criando ou removendo referências para inodes.
Elas são outro bom exemplo do poder do uso de 
transações. 
\lstinline{sys_link}
\lineref{kernel/sysfile.c:/^sys_link/}
começa obtendo seus argumentos, duas strings:
\lstinline{old}
e
\lstinline{new}
\lineref{kernel/sysfile.c:/argstr.*old.*new/}.
Assumindo que
\lstinline{old}
existe e não é um diretório
\linerefs{kernel/sysfile.c:/namei.old/,/^..}/},
\lstinline{sys_link}
incrementa seu contador
\lstinline{ip->nlink}.
Depois,
\lstinline{sys_link}
chama
\indexcode{nameiparent}
para encontrar o diretório pai e o último elemento do caminho de
\lstinline{new} 
\lineref{kernel/sysfile.c:/nameiparent.new/}
e cria uma nova entrada de diretório apontando para o inode de
\lstinline{old}
\lineref{kernel/sysfile.c:/\|\| dirlink/}.
O novo diretório pai deve existir e estar no mesmo dispositivo que o inode existente:
números de inode só têm significado único em um único disco.
Se um erro como esse ocorrer,
\indexcode{sys_link}
deve voltar e decrementar
\lstinline{ip->nlink}.

As transações simplificam a implementação porque é necessário atualizar múltiplos
blocos no disco, mas não precisamos nos preocupar com a ordem em que fazemos
essas atualizações. Ou todas terão sucesso ou nenhuma.
Por exemplo, sem transações, atualizar
\lstinline{ip->nlink}
antes de criar o link colocaria o sistema de arquivos temporariamente em um estado
inseguro, e uma falha no meio desse processo poderia resultar em corrupção.
Com transações, não precisamos nos preocupar com isso.

\lstinline{sys_link}
cria um novo nome para um inode existente.
A função
\indexcode{create}
\lineref{kernel/sysfile.c:/^create/}
cria um novo nome para um novo inode.
Ela é uma generalização das três chamadas de sistema para criação de arquivos:
\indexcode{open}
com o sinalizador
\indexcode{O_CREATE}
cria um novo arquivo comum,
\indexcode{mkdir}
cria um novo diretório,
e
\indexcode{mkdev}
cria um novo arquivo de dispositivo.
Como
\indexcode{sys_link},
\indexcode{create}
começa chamando
\indexcode{nameiparent}
para obter o inode do diretório pai.
Depois, chama
\indexcode{dirlookup}
para verificar se o nome já existe
\lineref{kernel/sysfile.c:/dirlookup.*[^=]=.0/}.
Se o nome já existir, o comportamento de 
\lstinline{create}
depende da chamada de sistema sendo usada:
\lstinline{open}
tem semânticas diferentes de 
\indexcode{mkdir}
e
\indexcode{mkdev}.
Se
\lstinline{create}
está sendo usado por
\lstinline{open}
(\lstinline{type}
\lstinline{==}
\indexcode{T_FILE})
e o nome existente for um arquivo regular,
então
\lstinline{open}
considera isso um sucesso,
e
\lstinline{create}
também
\lineref{kernel/sysfile.c:/^......return.ip/}.
Caso contrário, é um erro
\linerefs{kernel/sysfile.c:/^......return.ip/+1,/return.0/}.
Se o nome ainda não existir,
\lstinline{create}
aloca um novo inode com
\indexcode{ialloc}
\lineref{kernel/sysfile.c:/ialloc/}.
Se o novo inode for um diretório, 
\lstinline{create}
o inicializa com as entradas
\indexcode{.}
e
\indexcode{..}.
Por fim, agora que os dados estão devidamente inicializados,
\indexcode{create}
pode vinculá-lo ao diretório pai
\lineref{kernel/sysfile.c:/if.dirlink/}.
\lstinline{create},
como
\indexcode{sys_link},
mantém dois locks de inode simultaneamente:
\lstinline{ip}
e
\lstinline{dp}.
Não há possibilidade de deadlock porque
o inode
\lstinline{ip}
foi recém-alocado: nenhum outro processo no sistema
segurará o lock de
\lstinline{ip}
e depois tentará bloquear
\lstinline{dp}.

Usando
\lstinline{create},
é fácil implementar
\indexcode{sys_open},
\indexcode{sys_mkdir}
e
\indexcode{sys_mknod}.
\lstinline{sys_open}
\lineref{kernel/sysfile.c:/^sys_open/}
é o mais complexo, porque criar um novo arquivo é apenas
uma pequena parte do que ele pode fazer.
Se
\indexcode{open}
recebe o sinalizador
\indexcode{O_CREATE},
ele chama
\lstinline{create}
\lineref{kernel/sysfile.c:/create.*T_FILE/}.
Caso contrário, chama
\indexcode{namei}
\lineref{kernel/sysfile.c:/if..ip.=.namei.path/}.
\lstinline{create}
retorna um inode com lock, mas
\lstinline{namei}
não, então
\indexcode{sys_open}
deve fazer o lock do inode por conta própria.
Isso oferece um ponto conveniente para verificar que diretórios
só estão sendo abertos para leitura, não escrita.
Assumindo que o inode foi obtido de alguma forma,
\lstinline{sys_open}
aloca um arquivo e um descritor de arquivo
\lineref{kernel/sysfile.c:/filealloc.*fdalloc/}
e então preenche o arquivo
\linerefs{kernel/sysfile.c:/type.=.FD_INODE/,/writable/}.
Note que nenhum outro processo pode acessar o arquivo parcialmente inicializado,
já que ele está apenas na tabela do processo atual.

O Capítulo~\ref{CH:SCHED} examinou a implementação de pipes
antes mesmo de termos um sistema de arquivos.
A função
\indexcode{sys_pipe}
conecta essa implementação ao sistema de arquivos
fornecendo uma forma de criar um par de pipes.
Seu argumento é um ponteiro para um espaço de dois inteiros,
onde ela registrará os dois novos descritores de arquivo.
Depois, aloca o pipe e instala os descritores.

%% 
%%  -------------------------------------------
%% 
\section{Mundo real}

O cache de blocos em um sistema operacional real é significativamente
mais complexo que o do xv6, mas serve aos mesmos dois propósitos:
armazenar em cache e sincronizar o acesso ao disco.
O cache de blocos do xv6, como o do V6, usa uma política simples de substituição
de uso menos recente (LRU); há políticas muito mais complexas
que podem ser implementadas, cada uma adequada para alguns
tipos de carga e não tão boa para outros.
Um cache LRU mais eficiente eliminaria a lista encadeada,
usando em vez disso uma tabela de hash para buscas e uma heap para remoções LRU.
Caches modernos geralmente são integrados ao sistema de memória virtual
para suportar arquivos mapeados em memória.

O sistema de logging do xv6 é ineficiente.
Um commit não pode ocorrer simultaneamente com chamadas de sistema do sistema de arquivos.
O sistema registra blocos inteiros, mesmo que
apenas alguns bytes em um bloco sejam modificados. Ele realiza gravações
síncronas no log, um bloco de cada vez, cada uma provavelmente exigindo
uma rotação inteira do disco. Sistemas reais de logging tratam todos esses
problemas.

Logging não é a única forma de fornecer recuperação após falhas. Sistemas de arquivos antigos
usavam um varredor durante a reinicialização (por exemplo, o programa UNIX
\indexcode{fsck})
para examinar todos os arquivos e diretórios e as listas de blocos e inodes livres,
procurando e resolvendo inconsistências. Esse processo pode levar
horas em sistemas de arquivos grandes, e há situações em que não é
possível resolver inconsistências de forma que torne as chamadas de sistema originais atômicas.
A recuperação a partir de um log é muito mais rápida e garante que chamadas de sistema
sejam atômicas mesmo diante de falhas.

O xv6 usa o mesmo layout básico em disco de inodes e diretórios
que o UNIX antigo;
esse esquema tem se mantido notavelmente persistente ao longo dos anos.
O UFS/FFS do BSD e o ext2/ext3 do Linux usam essencialmente as mesmas estruturas de dados.
A parte mais ineficiente do layout do sistema de arquivos é o diretório,
que requer uma varredura linear sobre todos os blocos do disco em cada busca.
Isso é razoável quando os diretórios têm apenas alguns blocos,
mas é custoso para diretórios com muitos arquivos.
O NTFS do Windows, o HFS do macOS e o ZFS do Solaris, por exemplo,
implementam diretórios como árvores balanceadas em disco de blocos.
Isso é mais complexo, mas garante buscas em tempo logarítmico.

O xv6 é ingênuo quanto a falhas de disco: se uma operação de disco falha, o xv6 entra em pânico.
Se isso é aceitável ou não depende do hardware:
se o sistema operacional está rodando sobre hardware especial que usa
redundância para mascarar falhas de disco, talvez falhas sejam tão raras que a pane seja aceitável.
Por outro lado, sistemas operacionais que usam discos comuns
devem esperar falhas e lidar com elas de forma mais robusta,
de forma que a perda de um bloco em um arquivo não afete o
uso do restante do sistema de arquivos.

O xv6 requer que o sistema de arquivos
caiba em um único dispositivo de disco e que seu tamanho não varie.
À medida que grandes bancos de dados e arquivos multimídia aumentam
as necessidades de armazenamento, sistemas operacionais estão desenvolvendo formas
de eliminar o gargalo de ``um disco por sistema de arquivos''.
A abordagem básica é combinar vários discos em um único
disco lógico. Soluções de hardware como RAID ainda são as 
mais populares, mas a tendência atual é implementar
o máximo possível dessa lógica em software.
Essas implementações geralmente
permitem funcionalidades como expandir ou reduzir o dispositivo lógico dinamicamente.
Claro que uma camada de armazenamento que cresce ou diminui dinamicamente
requer um sistema de arquivos que possa fazer o mesmo: o array de tamanho fixo
de blocos de inode usado pelo xv6 não funcionaria bem
nesses ambientes.
Separar o gerenciamento de disco do sistema de arquivos pode ser
um design mais limpo, mas a interface complexa entre os dois
levou alguns sistemas, como o ZFS da Sun, a combiná-los.

O sistema de arquivos do xv6 carece de muitos recursos dos sistemas modernos; por exemplo,
não há suporte para snapshots e backups incrementais.

Sistemas Unix modernos permitem que muitos tipos de recursos sejam
acessados com as mesmas chamadas de sistema usadas para arquivos em disco:
pipes nomeados, conexões de rede,
sistemas de arquivos acessados remotamente e interfaces de monitoramento e controle
como o
\lstinline{/proc}.
Em vez das instruções
\lstinline{if}
em
\indexcode{fileread}
e
\indexcode{filewrite}
do xv6,
esses sistemas geralmente dão a cada arquivo aberto uma tabela de ponteiros de função,
um por operação,
e invocam o ponteiro correspondente à chamada implementada por aquele inode.
Sistemas de arquivos em rede e em espaço de usuário
implementam funções que transformam essas chamadas em RPCs de rede
e aguardam a resposta antes de retornar.

%% 
%%  -------------------------------------------
%% 
\section{Exercícios}

\begin{enumerate}

\item Por que fazer \lstinline{panic} em
\lstinline{balloc}?
O xv6 pode se recuperar?

\item Por que fazer \lstinline{panic} em
\lstinline{ialloc}?
O xv6 pode se recuperar?

\item Por que
\lstinline{filealloc}
não faz \lstinline{panic} quando os arquivos se esgotam?
Por que isso é mais comum e, portanto, vale a pena ser tratado?

\item Suponha que o arquivo correspondente a
\lstinline{ip}
seja desvinculado por outro processo
entre as chamadas de
\lstinline{sys_link}
para
\lstinline{iunlock(ip)}
e
\lstinline{dirlink}.
O link será criado corretamente?
Por que ou por que não?

\item
\lstinline{create}
faz quatro chamadas de função (uma para
\lstinline{ialloc}
e três para
\lstinline{dirlink})
das quais depende que todas tenham sucesso.
Se alguma falhar,
\lstinline{create}
chama
\lstinline{panic}.
Por que isso é aceitável?
Por que nenhuma dessas chamadas pode falhar?

\item
\lstinline{sys_chdir}
chama
\lstinline{iunlock(ip)}
antes de
\lstinline{iput(cp->cwd)},
o que pode tentar bloquear
\lstinline{cp->cwd},
mas adiar
\lstinline{iunlock(ip)}
até depois de
\lstinline{iput}
não causaria deadlocks.
Por que não?

\item Implemente a chamada de sistema
\lstinline{lseek}. Suportar
\lstinline{lseek}
também exigirá que você modifique
\lstinline{filewrite}
para preencher lacunas no arquivo com zeros se
\lstinline{lseek}
definir
\lstinline{off}
além de
\lstinline{f->ip->size.}

\item Adicione
\lstinline{O_TRUNC}
e
\lstinline{O_APPEND}
à
\lstinline{open},
de forma que os operadores
\lstinline{>}
e
\lstinline{>>}
funcionem no shell.

\item Modifique o sistema de arquivos para suportar links simbólicos.

\item Modifique o sistema de arquivos para suportar pipes nomeados.

\item Modifique o sistema de arquivos e o sistema de memória virtual para suportar arquivos mapeados em memória.

\end{enumerate}
